% Theory-owned Supplement S3 DRAFT; does NOT overwrite paper/5_selection.tex.
% Read together with S1 stop-time theory and registered v1/v2/v2-E1/v3 files.
\section{Residual selection, finite-panel estimators, and generation non-identifiability}
\label{sec:suppl-selection}

\paragraph{Class-conditional rank race.}
For one fixed image and solver, let $J\in\{G,O_g,O_p,C,I\}$
label a candidate with probabilities
$\theta=(p,o_g,o_p,c,u)$, and define $q=p+o_g+o_p$ and $s=q+c$.
The exchangeable-ranks reference assumes that, among candidates not
labeled I, residual rank does not favor a class. The observed first
2--4 runs contradict that calibration assumption: in the authorized
three-family A8 sample (1,960 reconstructions), GOOD outranks C in
200/205 paired within-image comparisons; for the 100 mixed-passing
image--solver records E0 selects GOOD 92 times, against 51.92 expected
under the exchangeable reference. The pair comparisons are dependent
and must not be treated as independent Binomial trials.

A conditional exponential-race \emph{surrogate for rank}, not a
physical-unit likelihood of the Fourier residual, assigns
$T\mid J=j\sim\operatorname{Exp}(w_j)$, independently across
runs, with $w_G=1$ and usually $w_I=0$. For a positive-rate class $j$,
\begin{equation}\label{eq:suppl-race-integral}
 W_j(K)=K\theta_j\int_0^\infty w_j e^{-w_jt}
 \Bigl(\sum_{\ell}\theta_\ell e^{-w_\ell t}\Bigr)^{K-1}dt.
\end{equation}
Conditioning on fixed class counts $N_j$ gives the competing-risks
identity
\[
 \Pr(\operatorname*{argmin}T\text{ has class }j\mid N)
 =\frac{N_jw_j}{\sum_{\ell}N_\ell w_\ell}
\]
whenever the denominator is positive. If the denominator is zero, a
predeclared tie convention is required. Passing to $K\to\infty$ at
fixed $\theta$ yields the selection ceiling
\[
 \lim_{K\to\infty}W_G(K)=
 \frac{p w_G}{p w_G+o_gw_O+o_pw_P+cw_C},
\]
provided the denominator is positive, instead of the reference
$p/s$. No generic monotonicity theorem applies to the oracle-minus-
selector gap under this race. Already with $p_G=p_{\rm bad}=1/2$ and
$(w_G,w_{\rm bad})=(1,.03)$, the gap obeys
$D_6=0.0386195271>D_7=0.0385396532$ and approaches
$D_\infty=0.0291262136$. The ordinary exponential race identity
and its asymptotics are classical; the solver-specific calibration
and prospective test constitute the present analysis.

\paragraph{Exact finite-panel and fixed-prefix estimands.}
Let the completed panel have $n$ runs in increasing physical-residual
rank order, with rank $r=1$ best. Uniformly choosing a $K$-subset
selects the candidate at rank $r$ with weight
\begin{equation}\label{eq:suppl-subset-weight}
 \omega_r(n,K)=\frac{\binom{n-r}{K-1}}{\binom nK}.
\end{equation}
The $K$-subset policy success is its exact weighted sum over the
candidate's actual GOOD label or orientation-corrected success
indicator. For an oracle that accepts $g$ candidates among the $n$,
\begin{equation}\label{eq:suppl-oracle-subset}
 \widehat O(K)=1-\frac{\binom{n-g}{K}}{\binom nK}.
\end{equation}
These are complete U-statistics, not independent fresh-$K$ trials.
The preregistered FFHQ/ImageNet forecasts hold original run indices
0--1 or 0--3 fixed and simulate *only* the remainder of the panel,
then evaluate the same finite-pool estimand. The additional optional
protocol theory must likewise condition on the original residual
threshold indicators and orientation decisions rather than assume
fresh unconditional negative-binomial panels.

\paragraph{Four-run count data do not identify a 64-run oracle.}
Let $p\in[0,1]$ be a latent per-image probability of returning raw GOOD,
and let its population mixing law be $\mu$. Four independent runs
produce $J\mid p\sim\operatorname{Bin}(4,p)$, so
\begin{equation}\label{eq:suppl-four-moments}
 \Pr(J=j)=\binom4j\int p^j(1-p)^{4-j}\mu(dp).
\end{equation}
These five probabilities determine only moments through degree four.
The following two measures are indistinguishable by *all* such
four-run count probabilities:
\[
 \mu_A=\operatorname{Uniform}[0,1],\qquad
 \mu_B=\frac19\delta_0+
 \frac{16+\sqrt6}{36}\delta_{(6-\sqrt6)/10}+
 \frac{16-\sqrt6}{36}\delta_{(6+\sqrt6)/10}.
\]
The classical three-node left-endpoint Gauss--Radau quadrature formula
makes the first four moments agree exactly. Yet the population oracle
$O_K=\int[1-(1-p)^K]\mu(dp)$ differs by
$O^A_{64}-O^B_{64}=0.095726\ldots$ and the limiting ceilings
differ by $1/9$. This is an *identification counterexample*, not a
model fitted to the FFHQ cohort or a proof that particular real
images have $p=0$.

\paragraph{Data-specific continuum outer bounds.}
An exploratory post-FAST64 calculation used only the historical A8
four-run count histograms for FAST to constrain the **population**
$O_{64}$. For each of the five count categories, take a marginal
99\% two-sided Clopper--Pearson confidence interval
$[L_j,U_j]$; Bonferroni yields simultaneous coverage at least 95\%
under iid image sampling. Search all probability measures $\mu$ with
\[
 L_j\le\int\binom4j p^j(1-p)^{4-j}\mu(dp)\le U_j
 \quad (j=0,\dots,4),\qquad \int\mu(dp)=1.
\]
With a grid of spacing $h=10^{-4}$, expand the count constraints by
$2h$ using the derivative bound $\|B_{4,j}'\|_\infty\le4$, and
expand the resulting minimum/maximum target by $32h$ using
$\|(1-(1-p)^{64})'\|_\infty\le64$. This gives conservative
\emph{continuum} outer enclosures, up to numerical solver tolerance,
\[
 O_{64}^{\mathrm{population}}\in
 \begin{cases}
 [0.8696,1],&\text{FFHQ cohort A (100 images)},\\
 [0.6113,1],&\text{FFHQ enriched cohort B (80 images)}.
 \end{cases}
\]
This bound is \emph{not} a confidence interval for the frozen
fixed-image conditional-prefix statistic. The nominal frequentist
coverage requires iid image sampling within the cohort; enrichment
of cohort B requires an additional assumption, otherwise its
interval is descriptive sensitivity. The LP was designed after
seeing the FAST64 miss, notwithstanding its use of historical
four-run counts only: it is not a preregistered prediction.

\paragraph{Interpretation of the two prospective tests.}
On FAST FFHQ, the preregistered E0 rank-race interval covered 6/21
primary-panel points against 15/21 for the exchangeable-ranks
reference, while both models overforecast the oracles. This is an
\emph{unconditional generation failure}, not independent evidence
in favor of exchangeable ranks: exploratory conditioning on the
observed class composition gives pooled K=64 race E0 0.9526 vs
0.9389 observed, while exchangeable ranks give 0.8773.
On ImageNet D1, the separately preregistered orientation-corrected
E1 addendum covered 6/6 checks for the race variants and 0/6 for
S2. At K=16, E1 was 0.650 vs race 0.603 [0.520,0.682], but the
orbit oracle was only 0.750 vs 0.859 forecast (outside its
predictive interval). These are different decompositions and
should be reported jointly: prediction of the deployed E1 output
does not validate all of its generative components. Neither held-out
experiment justifies retrospective refitting of its registered model.
